\documentclass[main.tex]{subfiles}


\begin{document}

% \AddToShipoutPictureFG{%
% \begin{tikzpicture}[overlay,remember picture]
% \path (current page.south west) -- (current page.north east)
% node[midway,scale=1,color=gray,sloped,opacity=0.4] {\textnormal{Кравченко Игорь Игоревич\hspace{1cm} +79010144910\hspace{1cm} super.antig@yandex.ru}};
% \end{tikzpicture}
% }

% %from pagenumber to up
% \phantomsection 
% \hypertarget{toc}{}

 

 

\setcounter{section}{1}
\addtocounter{section}{-1} 

\section{Векторы}


\textbf{Вектор}~--- это стрелка на чертеже для описания направления физической характеристики. На рис.\,\ref{fig:p1}
вектор имеет синий цвет.

\begin{figure}[H]
    \centering

    \begin{tikzpicture}[font=\small,            line cap=round,                             >={Stealth[inset=0pt,round,length=3mm, angle'=20]},vec/.style={>={Stealth[inset=0pt,round,length=3.5mm, angle'=20]}}]


\draw[->,NavyBlue,thick,vec] (0,0) --++ (0:2cm) node[midway,above,black] {$\vec v$};
\draw[red,arrows = {Bar[line cap=round]-Bar[line cap=round]}] (5,0) --++ (0:1cm) node[midway,above,black] {3~м$/$с};
\path[|-|,red,thick,vec] (5,0) --++ (0:1cm) node[midway,below,black] {$m_v$};


    \end{tikzpicture}
\caption{Вектор и масштаб}
\label{fig:p1}
\end{figure}

% \begin{figure}[H]
%     \centering

%     \includestandalone{PICS/mech/1vec_1}

% \caption{Вектор и масштаб}
% \label{fig:p1}
% \end{figure}



Красный элемент на рис.\,\ref{fig:p1} есть \emph{масштаб}, по которому определяется длина вектора в указанных единицах~--- то есть \textbf{величина} или \textbf{модуль вектора}. Видно, что величина данного вектора $v=6$~м$/$c. Используются также и другие единицы измерения (буквы после числа)~--- например, <<м>>, <<м$/$с$^2$>> и~т.\,д. 

\textbf{Проекция вектора}~--- это <<тень>> вектора на ось координат при освещении перпендикулярно на эту ось. На рис.\,\ref{fig:p2} показаны проекции $v_x$ и $v_y$. 

\begin{figure}[H]
    \centering

    \begin{tikzpicture}[font=\small,            line cap=round,                             >={Stealth[inset=0pt,round,length=3mm, angle'=20]},vec/.style={>={Stealth[inset=0pt,round,length=3.5mm, angle'=20]}}]


%light and stuff
\filldraw[yellow] (0.5,3) -- (3.5,3) -- (3.5,0) -- (0.5,0) -- cycle;
\filldraw[white] (1,1) -- (3,2) -- (3,0) -- (1,0) -- cycle;
\draw[->,orange] (4,3) --++ (0,-1) node[midway,right,black] {Свет на $Ox$};


%xoy
\draw[->] (-.3,0) --++ (0:5cm) node[below,black] {$x$};
\draw[->] (0,-.3) --++ (90:3.3cm) node[left,black] {$y$};
\node[below left] at (0,0) {$O$};

%vec
\draw[->,thick,vec,NavyBlue] (1,1) -- (3,2) node[midway,below,black] {$\vec v$};

%pro
\path[-,thick,vec,lightgray,decoration={markings,mark=at position 1 with {\arrow{>}}},postaction={decorate}] (1,0) -- (3,0);
\draw[-,ultra thick,vec,gray] (1,0) -- (3,0) node[midway,below,black] {$v_x$};

\path[-,thick,vec,lightgray,decoration={markings,mark=at position 1 with {\arrow{>}}},postaction={decorate}] (0,1) -- (0,2);
\draw[-,ultra thick,vec,gray] (0,1) -- (0,2) node[midway,left,black] {$v_y$};

%dashed
\draw[densely dashed] (1,1) -- (0,1); 
\draw[densely dashed] (3,2) -- (0,2); 

    \end{tikzpicture}
\caption{Проекции вектора}
\label{fig:p2}
\end{figure}

% \begin{figure}[H]
%     \centering

%     \includestandalone{PICS/mech/1vec_2}

% \caption{Проекции вектора}
% \label{fig:p2}
% \end{figure}


В общем случае проекция вектора равна длине <<тени>> со знаком
\begin{itemize}
    \item $+$, если <<теневая>> стрелка стоит по оси;
    \item $-$, если <<теневая>> стрелка стоит против оси.
\end{itemize}

На рис.\,\ref{fig:p2} по обозначениям возле <<теней>> видно, что проекция не является вектором. Если приписать <<крыши>> над $v_x$ и $v_y$ и получше выделить их <<теневые наконечники>>, то окажется, что вектор~$\vec v$ разложен на \textbf{составляющие} векторы~$\vec v_x$ и~$\vec v_y$. Треугольник из этих векторов дает: $v=\sqrt{\mathstrut v_x^2+v_y^2}$.   

\textbf{Сложить векторы}~--- это значит найти вектор, который совмещает в себе как бы <<действия>> этих векторов. Два сонаправленных вектора <<дают>> вектор в ту же сторону (длины складываются). Также ясно, что два противоположно направленных вектора <<дают>> вектор в сторону большего  вектора (длины вычитаются). На рис.\,\ref{fig:p3} показано, как сложить два произвольных вектора.  

\begin{figure}[H]
    \centering

    \begin{tikzpicture}[font=\small,            line cap=round,                             >={Stealth[inset=0pt,round,length=3mm, angle'=20]},vec/.style={>={Stealth[inset=0pt,round,length=3.5mm, angle'=20]}}]

%vec
\draw[->,NavyBlue,thick,vec] (0,0) --++ (60:2cm) node[midway,above left,black,inner xsep=0] {$\vec v_1$};
\draw[->,NavyBlue,thick,vec] (1,0) --++ (-40:1cm) node[midway,below left,black,inner xsep=0] {$\vec v_2$};

\draw[->,NavyBlue,thick,vec] (1,0) ++ ({cos(-40)},0) ++ (4,0) node (12) {} --++ (60:2cm) node[midway,above left,black,inner xsep=0] {$\vec v_1$};
\draw[->,NavyBlue,thick,vec] (12.center) --++ (-40:1cm) node[midway,below left,black,inner xsep=0] {$\vec v_2$};


%dash
\draw[dashed,gray,shift={(-40:1cm)}] (1,0) ++ ({cos(-40)},0) ++ (4,0) --++ (60:2cm);
\draw[dashed,gray] ([shift={(60:2cm)}]12.center) --++ (-40:1cm) node (n) {};


%vec12
\draw[->,red,thick,vec] (12.center) -- (n.center) node[right,black] {$\vec v_{12}$};

%\Longrightarrow
\path[->,thick,vec] ({1+cos(-40)},{(sin(60)*2-sin(40))/2})  --++ (4,0) node[midway] {$\Longrightarrow$};

    \end{tikzpicture}
\caption{Правило параллелограмма}
\label{fig:p3}
\end{figure}

% \begin{figure}[H]
%     \centering

%     \includestandalone{PICS/mech/1vec_3}

% \caption{Правило параллелограмма}
% \label{fig:p3}
% \end{figure}



 
\end{document}